\head{\textsc{\emph{Ilíada}, Canto V}}
\section*{\centering\textbf{\textsc{Canto} V}\footnote{\ehdos, pp. 113-131.}}
\addcontentsline{toc}{subsection}{\textsc{Canto V}}
%
\beginnumbering
\pstart
\firstlinenum{5}
\lettrine[lines=4,nindent=2pt]{P}{\textsc{alas Atena, entonces}}, dió audacia y entereza\\
a Diomedes Tideides, para que alta proeza\\
destacara entre todos los argivos su fama.\\
Le encendió escudo y yelmo con incansable llama,\\
tal como la del astro de Otoño, que lavado\\
por el Océano emerge más que todos brillante:\\
así el fuego en sus hombros y cabeza ha brotado,\\
cuando lo echa en lo recio del tumulto, adelante.\par
Había en Troya un cierto Dares, rico y honrado,\\
sacerdote de Hefesto, cuyos dos hijos, que eran\\
Ideo y Fegeo, aptos para todo combate,\\
separados del resto, contra el héroe se abaten\\
montando un carro, mientras él a pie los espera.\\
Fegeo, al encontrarse , tiró primeramente\\
su venablo de larga sombra que inútilmente\\
por sobre el hombro izquierdo del Tideida pasó.\\
Pero éste con segura mano le arrojó el bronce\\
que en la mitad del pecho fué a clavársele entonces,\\
tumbándolo del carro, pues no en vano partió.\\
Corrió Ideo, el hermoso vehículo dejando,\\
sin defender el cuerpo del hermano caído.\\
Que él mismo de la negra muerte no hubiera huído,\\
si en tinieblas Hefesto no lo envuelve, evitando\\
con salvarlo, que el viejo quede más afligido,\\
y el hijo de Tideo tomó y dió los corceles\\
a los suyos, haciéndolos llevar a los bajeles.\\
Al notar los troyanos que huye un hijo de Dares\\
y que el otro ante el carro muerto quedó, se apena\\
su corazón magnánimo; y la ojizarca Atena\\
así dijo, domándole la mano, al furioso Ares:\par
--Ares, Ares, funesto, destructor, sitiador:\\
¿no es mejor que dejemos que aqueos y troyanos\\
luchen, y que la gloria les dé Zeus soberano,\\
mientras así eludimos su divino rencor?\par
Tal dijo, y del combate sacando a Ares furioso,\\
lo hizo sentar a orillas del Escamandro herboso.\\
Los dánaos pusieron en fuga a los troyanos,\\
y cada jefe, entonces, mató un hombre a sus manos.\\
Agamenón, rey de hombres, va primero y derriba\\
del carro al grande Odíon, rey de los halizones,\\
que dió la espalda, hundiéndole su lanza desde arriba,\\
y el pecho atravesándole por entre los pulmones.\\
Tal cayó, y con estrépito sus armas resonaron.\par
Mató allá Idomeneo, por la lanza preclaro,\\
a Festo, hijo de Boro el meonio, que había\\
llegado de la fértil Tarne, hundiendo la enorme\\
asta en su hombro derecho, cuando al carro subía.\\
Cayó así; horribles sombras entonces lo rodearon,\\
y los hombres de Idomeneo lo despojaron.\par
Menelao el Atrida mató a Escamandrio, experto\\
cazador, que de Estrofio fué hijo, con su asta de haya.\\
La propia Artemis era quien le enseñó el acierto\\
en cazar cuantas fieras en bosques y cerros haya.\\
Mas de nada valiéronle ni artemis la flechera,\\
ni su celebrado arte de tirar al acecho,\\
pues el ilustre Atrida que huyendo ante él lo viera,\\
lo lanceó entre los hombros, traspasándole el pecho.\\
Y así, al caer de bruces, sus armas resonaron.\\
--Merión mató a Fereclón, hijo de Harmón maestro\\
que fuera en toda clase de artes manuales diestro,\\
pues Palas Atenea le daba su alto amparo.\\
Fué él quien hizo a Alejandro los navíos veloces\\
que tanto mal a Troya y a él mismo ocasionaron,\\
por no haber entendido su oráculo a los dioses.\\
Merión al alcanzarlo, pues sin cesar lo hostiga,\\
por la nalga derecha le envasó el bronce fuerte;\\
y como el hueso hendiérale, tocando la vejiga,\\
gimió al caer de hinojos y lo envolvió la muerte.\par
Meges mató a Pedeo, bastardo de Antenor,\\
a quien como hijo propio la divina Teano,\\
por ser grata a su esposo, criara con amor.\\
El ilustre Fileides que a él se hallaba cercano,\\
fué y le clavó en la nuca su aguda lanza entonces,\\
y la lengua cortándole, le atravesó los dientes:\\
que así rodó por tierra, mordiendo el frío bronce.\par
Eurípilo Evemónides da al divino Hipsenor\\
hijo de Dolopión el Grande, que oficiaba\\
el culdo de Escamandro, y a quien el pueblo honraba\\
como a un dios, fiero tajo cuando huye con pavor.\\
La espada, junto al hombro, tronchó el brazo pesado,\\
y rodó por los suelos el miembro ensangrentado.\\
Y cerraron los ojos de aquél en la cruenta\\
lid, la muerte purúrea y la Parca violenta.\par
Y mientras así agítalos el combate empeñoso,\\
no se sabe si el hijo de Tideo, en su brío,\\
combate con troyanos o aqueos, tan furioso\\
lánzase a la llanura cual rebalsado río\\
que durante el deshielo va arrollando bravío\\
los diques y los setos de los campos feraces,\\
y al crecer con las lluvias de Zeus, en su desvío,\\
muchas bellas labranzas de los mozos deshace:\\
tal las densas falanges troyanas descompone\\
Tideides, y entre tantos, ninguno se le opone.\\
Cuando, en eso, el caro hijo de Licaón, que advierte\\
que en la tropa del llano siembra, feroz, la muerte,\\
tiende el corvo arco y clávale en la espalda derecha,\\
por el hueco de la honda coraza, amarga flecha\\
que ensagrentando el bronce, del otro lado asoma.\\
Y el claro Licaónida grita con voz terrible:\par
--¡Cargad, bravos troyanos, hábiles en la doma!\\
Herido está el aqueo más bravo, y no es posible\\
que el grave tiro aguante, si con suerte propicia\\
me endilgó acá el rey, vástago de Zeus, desde la Licia.\par
Tal se alaba, aunque al otro no rindió el arma aguda;\\
pues ante yunta y carro reculó con pericia,\\
pidiendo a Esténelo, hijo de Capaneo ayuda:\par
--¡Corre, buen Capaneides, deja el carro, y prolijo\\
\dlgo{arráncame del hombro la amarga flecha!}{Dijo,}\\
y al punto apeóse Esténelo del carro, y la enemiga\\
flechada clavada al hombro le arrancó; y desbordando,\\
saltó la sangre por las mallas de la loriga.\\
Y así fue el buen guerrero Diomedes implorando:\par
--Escúchame hija indómita de Zeus potente: si\\
tu amparo en la ardua lucha nos diste siempre a mí\\
y a mi padre, ahora haciéndome la gracia, oh Atenea,\\
de que mate yo a ese hombre, ponlo a tiro de lanza;\\
pues porque me hirió, alábase con la loca esperanza\\
de que yo, pronto, el caro fulgor del sol no vea.\par
Así rogó y fué oído por Palas Atenea,\\
qe aligerando al punto su cuerpo, pies y manos,\\
con aladas palabras lo animó a la pelea:\par
--Recóbrate, Diomedes, y ataca a los troyanos,\\
que ya infundí en tu pecho la audacia y el vigor\\
de tu intrépido padre Tideo el domador,\\
hábil en el escudo; y aun torné caediza\\
la nube de ceguera que tapaba tus ojos,\\
para que a dioses y hombres conozcas en la liza,\\
y si un numen te prueba, no afrontes sus enojos.\\
Sólo entre Afrodita, la hija de Zeus, entonces\\
hiérela sin escrúpulo con el agudo bronce.\par
Dicho esto, fuése Atene la ojizarca; y al frente\\
de combate, Tideides se incorporó sin falta,\\
pues contra los troyanos, más aún que anteriormente,\\
triplicósele el brío, como al león que asalta\\
un establo de ovejas, donde lo ha lastimado\\
el pastor sin rendirlo, con que su furia exalta.\\
Húyese el hombre, entonces, al fondo del cercado;\\
y mientras indefenso y en su terror febril\\
con revuelto atropello se amontona el ganado,\\
la ávida fiera salta del profundo redil.\\
Así el fuerte Diomedes, en furia arrebatado,\\
las falanges troyanas revuelve y desbarata.\par
Allá a Astino y a Hipéiron, pastor de pueblos, mata,\\
clavando a éste la pica de bronce en la tetilla,\\
mientras con la ancha espada troncha al otro la istilla,\\
y de la espalda y cuello todo el hombro separa.\\
Déjalos allá y cierra contra Abante y Polido\\
a los que un viejo intérprete de sueños engendrara:\\
Eurídamas, que al verlos partir no habrá entendido\\
los suyos, pues Diomedes a ambos mata y despoja.\\
Acto continuo contra Janto y Toón se arroja,\\
hijos que pudo Fénope en su vejez lograr,\\
pues ya caduco y triste, no engendró otro heredero.\\
Al quitarles la amable vida, en llanto y pesar\\
sumió al padre, que vivos no ha de verlos tornar:\\
con los que otros parientes la herencia se repartieron.\par
Luego en un carro toma dos héroes que hijos fueron\\
de Príamo Dardánida: Cromíon y Equemón.\\
Cual desnuca, asaltando la boyada, el león,\\
a un buey o una ternera que en el bosque pacía,\\
Tideidos lo desmonta por la fuerza, con impía\\
saña, pilla sus armas, y entrega los corceles\\
a los suyos, haciéndolos llevar a los bajeles.\par
Mas, cuando destrozaba las filas, viólo Eneas,\\
y por entre el tumulto de armas de la pelea\\
fué en busca del divino Pándaro. Así que halló\\
al eximio y bravo hijo de Licáon, lo encaró\\
\dlgo{para decirle:}{--Pándaro, ¿qué es de tu arco y aladas}\\
flecha, y de tu fama por nadie disputada\\
aquí ni en Licia, donde ninguno se gloría\\
de aventajarte? ¡Vamos! levanta a Zeus tus manos,\\
y tira contra ese hombre, que haciendo a los troyanos\\
innumerables males, triunfa con saña impía,\\
y así a tanto valiente las rodillas afloja;\\
si no es algún dios que hacia los de Troya indispuso\\
un sacrilegio, porque malo es si un dios se enoja.\par
Entonces el claro hijo de Licáon le repuso:\\
--Eneas que a los nobles troyanos aconsejas,\\
en todo al bélico hijo de Tideo semeja.\\
Su escudo, su oculario triple casco ver creo,\\
y sus corceles, aunque puede que sea un dios;\\
mas, si es el bélico hijo del ilustre Tideo,\\
su furia está probando que un inmortal va en pos,\\
envuelto en una nube, y apartando la aguda\\
flecha que lo amenaza; pues ya una le tiré,\\
y por el hueco de la coraza le acerté\\
a la espalda derecha, creyendo que sin duda\\
al Hades lo echaría, pero no lo maté:\\
con que, seguramente, será un dios irritado.\\
Ni caballos ni carro tengo aquí qué montar,\\
aun cuando en el palacio de Licáon han quedado\\
once hermosos y sólidos coches sin estrenar,\\
con los toldos dispuestos, y cada uno a su lado\\
la yunta que cebada blanca y avena come.\\
El anciano guerrero Licáon prescribióme\\
al dejar yo el palacio, con premiosa insistencia,\\
que a los troyanos sólo desde el carro mandara\\
en la ardua lucha, pero no le presté aquiesencia.\\
Cuánto más me valiera, pues dejé por prudencia\\
los caballos, temiendo que el pasto les faltara\\
con reunión tan grande, y estando acostumbrados\\
a comer bien; sin ellos, a pie a Ilión he venido,\\
confiado sólo en mi arco, que nada útil me ha sido;\\
pues ya contra dos jefes, Tideides y un Atrida\\
he tirado, y aunque a ambos sangré con grave herida,\\
sólo excitarlos pude. Con mal sino partí\\
el día en que el corvo arco del clavo desprendí\\
y a la amena Ilión traje mis troyanos, por ser\\
amable al divino Héctor. Mas, como vuelva a ver\\
a mi patria, a mi esposa y a mi alta fortaleza,\\
que un guerrero enemigo me corte la cabeza,\\
si este arco no destrozo y echo al fuego violento,\\
ya que su compañía me es vana comp el viento.\par
Mas, replicóle el jefe de troyanos, Eneas:\\
--No hables así, pues nada cambiará, si nosotros\\
no vamos contra ese hombre con armas, carro y potros\\
a pelearlo. ¡Ea! al mío sube, para que veas\\
cómo saben lo mismo los corceles troyanos\\
perseguir o escaparse rápido por los llanos,\\
y cómo han de llevarnos salvos a la ciudad,\\
si otra vez Zeus concede gloria, en su potestad,\\
a Diomedes Tideides. Toma, pues, ya, a mi lado,\\
fusta y brillantes riendas; yo pelearé montado,\\
o hazlo tú, y los caballos a mi cuidado sean.\par
\dlgo{Y contestó el claro hijo de Licaón:}{--Eneas,}\\
encárgate tú mismo de la yunta y las bridas;\\
pues hallándose bajo su habitual conductor,\\
si el hijo de Tideo nos forzara a la huída,\\
tirará aquélla el cóncavo carro mucho mejor;\\
no se asuste y desboque, tu voz desconociendo,\\
y a sacarnos se niegue de la lucha, y cayendo\\
sobre ambos el bravo hijo de Tideo, nos mate\\
y se lleve los fuertes potros. Así, al combate\\
conduce por tu propia mano tu carro y yunta,\\
mientras su ataque espero con mi lanza de punta.\par
Hablando así, al labrado carro suben, y ansiosos,\\
contra Tideides lanzan los corceles briosos.\\
Mas Esténelo, el célebre hijo de Capaneo,\\
que los advierte, dícele con palabras aladas:\par
--Carísimo Diomedes Tideides, venir veo\\
dos robustos varones, con ardiente deseo\\
de pelearte, y que tienen fuerzas desmesuradas:\\
Pándaro el diestro arquero, quien se jacta de que es\\
Vástago de Licáon, y Eneas, que a su vez\\
se alaba de que Anquises, varón irreprochable,\\
lo engendró en Afrodita. Retirémonos, pues,\\
montados; no te expongas, furioso, a algún revés,\\
si en la vanguardia arriesgas perder la vida amable.\par
Mas el fuerte Diomedes, mirándolo de abajo,\\
\dlgo{contestó:}{--No me digas que huya o te des trabajo}\\
de convencerme en vano, pues no soy del linaje\\
de los que en retirada, temerosos pelean.\\
Iré a ellos desmontado, que aun me sobra coraje,\\
firme con el apoyo de Palas Atenea.\\
No han de llevarlos lejos de acá, si retroceden,\\
sus veloces caballos, y alguno escapar puede.\\
Mas, te diré algo, y guárdalo fielmente en tu memoria:\\
si la sapiente Atena me concede la gloria\\
de matarlos, tú nuestros rápidos troncos ata\\
al barandal, ciñéndoles las riendas, y arrebata\\
lejos de los troyanos los corceles de Eneas,\\
que pondrás al amparo de las tropas aqueas;\\
pues su raza es la de esos que a Tros dió en expiación\\
por Ganimedes, su hijo, Zeus el clarividente.\\
Nada hay, bajo la aurora y el sol, más excelente.\\
El rey de hombres Anquises, burló a Laomedón\\
y sacó cría, echándoles yeguas ocultamente.\\
Naciéronle seis de ese linaje en su mansión:\\
cuatro crió a pesebre, y a Eneas hizo don\\
de los otros dos, que ahora sembrar el terror vemos.\\
Con que, si los tomamos, un buen triunfo obtendremos.\par
Así éstos conversaban, mientras que los dos otros\\
caían ya sobre ellos con sus veloces potros.\\
Y este reto el claro hijo de Licáon les echa:\par
--¡Corajudo y noble hijo del ilustre Tideo,\\
pues no llegó a postrarte la amarga y rauda flecha,\\
veremos si la lanza responde a mi deseo!\par
Dice, y cimbrando tira su asta de larga sombra,\\
y el broquel del Tideida con el bote atraviesa,\\
rozando la coraza, y así a jactarse empieza:\par
--¡Atravesado tienes el vientre, y no me asombra\\
que pronto acabes, dándome considerable fama!\par
Pero el fuerte Diomedes, sin conmoverse, exclama:\\
--¡Erraste y falló el tiro! Mas creo por mi parte\\
que no lograréis tregua, como uno de ambos no harte\\
de sangre al bélico Ares, rodando por la arena.\par
Dice y tírale; el dardo, guiado por Atena,\\
entre ojo y nariz clávasele, sus blancos dientes pasa,\\
y el incansable bronce la lengua le cercena\\
de raíz, y su punta bajo el mentón rebasa.\\
Así cayó del carro; resonaron sobre él\\
sus magníficas armas, se abrieron sus corceles,\\
y allá lo abandonaron su fuerza y su alma fiel.\par
Eneas salta al suelo con su lanzón y escudo,\\
y temiendo que el cuerpo le quiten los aqueos,\\
ronda en torno, confiado como un león forzudo,\\
prontos broquel y pica, con ardientes deseos\\
de matar al que ataque, mientras grita sañudo.\\
Tideides alza entonces una piedra imponente\\
que dos hombres de ahora levantar no podrían,\\
mientras con una mano la lleva él fácilmente.\\
Con ella hiere a Eneas donde el fémur encaja\\
con la cadera, punto que se llama \emph{cotylo};\\
rómpele ambos tendones y el \emph{cotylo} descuaja,\\
y el cutis le desgarra la piedra con su filo.\\
El héroe, de rodillas cae, apoyando en tierra\\
su fuerte mano, y negra noche sus ojos cierra.\par
Y allá Eneas, rey de hombres, habría perecido,\\
si acaso Afrodita, hija de Zeus, que lo tuviera\\
de Anquises al boyero, pronto no lo advirtiera.\\
Con que, sus blancos brazos tiende al hijo querido,\\
y en un doblez del mando brillante lo resguarda,\\
no sea que los dánaos bien montados le claven\\
el cruel bronce en el pecho, y así su vida acaben.\par
Y mientras del combate lo substrae, no tarda\\
en efectuar Esténelo lo que convino, al mando\\
del cumplido guerrero Diomedes. Apartando\\
la yunta del tumulto, sus bridas ciñe al tope\\
del barandal, y leos de los troyanos lleva\\
los lucientes caballos de Eneas, al galope,\\
rumbo hacia los aqueos de las hermosas grebas\\
los da a Deipilo (amigo que entre todos honraba,\\
por lo bien que su juicio con el suyo acordaba)\\
a fin de que a los buques los conduzca; y montando\\
su yunta otra vez, coge más bien labradas riendas,\\
y solícito anímala, a Tideides buscando.\\
Este acosa a Ciprina con sus armas tremendas,\\
sabiéndola una débil diosa entre las deidades,\\
no de esas que a los hombres guían en las contiendas,\\
como Atena o Enío que arrasa las ciudades.\\
Y así que en la refriega tumultuosa la alcanza,\\
el hijo de Tideo cala su aguda lanza\\
y con ella atacándola, rompe el mando divino,\\
propia obra de las Gracias, y herir el cutis fino\\
de la mano, al extremo del dorso palmar, osa.\\
Brota entonces la sangre divina de la diosa,\\
o mejor dicho el ícor, flúido peregrino\\
que los dioses dichosos en vez de ella derraman,\\
porque ni trigo comen, ni beben negro vino:\\
(con lo que, siendo exangües, inmortales los llaman).\\
Lanza ella agudo grito, dejando a su hijo solo;\\
y mientras lo alza en una nube azul Febo Apolo,\\
no sea que los dánaos bien montados le claven\\
el cruel bronce en el pecho, y así su vida acaben,\\
increpa el buen guerrero Diomedes a la diosa:\par
--¡Hija de Zeus, apártate de la acción belicosa!\\
¡Basta ya con que engañes a débiles mujeres!\\
Pues creo que si en otra guerra mezclarte quieres,\\
hasta su mismo nombre te alejará medrosa.\par
Dijo, y ella apartóse cruelmente torturada.\\
La rauda Iris sacóla del tumulto, abrumada,\\
negreándole ya el lindo cutis. Pronto halló al fiero\\
Ares, que hacia la izquierda del combate sentado,\\
contra una nube había su lanzón arrimado,\\
lo propio que la yunta de caballos ligeros.\\
Y ante el querido hermano se arrodilló, abatida,\\
instándole a prestarle los potros de áurea brida:\par
--Querido hermano, apiádate de mi y préstame ahora,\\
para ir al Olimpo en que los inmortales moran,\\
tus caballos; pues sufro mucho de hallarme herida\\
por un varón, Tideides, quien con Dios padre fuera\\
\dlgo{capaz de Pelear.}{Díjole así, y Ares le ha dado}\\
los potros de áurea brida, con que el carro ha montado,\\
gimiendo de congoja. Subió Iris a su vera,\\
y picó, rienda en mano, la yunta voladora.\\
Pronto al Olimpo fueron, donde los dioses moran;\\
y la diligente Iris, los potros sujetando,\\
los desató y echóles por pasto la ambrosía.\par
La divina Afrodita, mientras tanto, caía\\
al regazo de Dione su madre, que tomando\\
en los brazos a su hija, la acarició y le dijo:\par
--¿Qué uraniano, hija amada, te hizo esta iniquidad,\\
como si procedieras con notoria maldad?\par
\dlgo{La risueña Afrodita respondióle:}{--Fué el hijo}\\
de Tideo, el soberbio Diomedes, quien me hirió,\\
porque a Eneas, que es mi hijo más caro, saqué yo\\
del combate. Pues ahora ya no son entre iguales,\\
aqueos y troyanos, las luchas más furiosas;\\
que los dánaos atacan hasta a los inmortales.\par
Y le contestó Dione, divina entre las diosas:\par
--Soporta, hija, tus penas, que muchos moradores\\
del Olimpo, debimos sufrir tormento agudo\\
de los hombres, mezclándonos así con sus rencores.\\
Los pasó Ares, cuando Oto y Efialtes el forzudo,\\
hijos de Alo, apresáronlo con formidable nudo,\\
por trece meses dentro de un cántaro de bronce.\\
Y muriera allá el bélico Ares, si Eríbea, entonces,\\
bellísima madrastra de aquéllos, no lo advierte\\
a Hermes, quien substrajo a Ares del lazo rudo y fuerte\\
que lo postraba exánime. Hera los padeció,\\
cuando el vigoroso hijo de Anfitrión le clavó\\
en el seno derecho saeta trifurcada,\\
en tormento incurable dejándola angustiada.\\
Y sufrió el enorme Hades la fecha audaz de aquel\\
mismo hombre, hijo de Zeus Portaégida, cuando ante\\
la puerta de los muertos le infligió herida cruel.\\
De dolor traspasado, y el corazón penante,\\
pues abatía su ánimo la saeta enterrada\\
en su espalda robusta, fué al Olimpo, morada\\
de Zeus, donde aplicándole Peón drogas calmantes,\\
lo curó, pues no había nada en él de mortal.\\
¡Bárbaro y temerario quien con impiedad tal,\\
por el arco a los dioses del Olimpo abatió!\\
Así Atena, la zarca diosa, contra tí alzó\\
al hijo de Tideides, quien, insensato, ignora\\
que todo el que a los dioses combate en mala hora,\\
no alcanza larga vida, ni logra que ocupando\\
sus rodillas los hijos, le llamen padre, cuando\\
vuelve de los combates y la feroz pelea.\\
Piense Tideides, aunque se tenga por tan fuerte,\\
que alguien puede atacarlo mejor que tú; no sea\\
que muy luego la próvida Adrastina Egialea,\\
corte, gimiendo, el sueño, y en su inquietud despierte\\
a las caras doncellas, llorando a su marido\\
legítimo, ese aqueo principal y cumplido,\\
el domador Diomedes, de quien es noble esposa.\par
Dice y resta el ícor, en las suyas tomando\\
la mano que así cura, su cruel dolor calmando.\\
Mas, como Atena y Hera presenciaran la cosa,\\
picaron con palabra mordaz al dios Kronida.\\
Y dijo, adelantándose, así la zarca diosa:\par
--Zeus padre, aunque te irrite mi palabra atrevida,\\
seguro es que Ciprina, tentando a alguna aquea\\
de hermoso peplo, para que siga a los troyanos\\
que tanto ama, al hacerle las caricias que emplea,\\
se arañó en áureo bronce la delicada mano.\par
Tal dijo. El padre de hombres y dioses sonrió,\\
y a la áurea Afrodita llamando, así le habló:\par
--El bélico trabajo no es para tí, hija mía,\\
sino los de himeneo, con su dulce alegría.\\
Que cuiden de eso Atena y Ares el dios ligero.\par
Mientras así conversan, lánzase el buen guerrero\\
Diomedes contra Eneas, aunque sabe que Apolo\\
le tiene de su mano, pues ni al gran dios flechero\\
respeta ya, deseando matar a aquél tan sólo,\\
y quitarle el magnífico arnés. Tres veces carga\\
con ansia de ultimarlo, y otras tantas Apolo\\
se le opone, agitando su refulgente adarga.\\
Mas, cuando igual a un numen por cuarta vez se obstina,\\
Apolo, el grande arquero, terrible lo amenaza:\par
--Tidida, piensa y vete; no quieras darte traza\\
de igualarte a los dioses, que la raza divina\\
nunca se asemejó a la terrestre humana raza.\par
Así dice, y Tideides recula un poco, sólo\\
por evitar la cólera del flechador Apolo.\par
\pend
%
\noindent-- -- -- -- -- --\par
%
\pstart
\setline{584}Antíloco, atacándolo, le hundió en la sien la espada,\\
con que, del carro hermoso, cayó al suelo, anhelante,\\
quedando un largo rato de cabeza, clavada\\
toda ella hasta los hombros en la arena abundante.\\
Tendiólo al fin por tierra la yunta alborotada,\\
que picó hacia el ejército aqueo el valeroso\\
Antíloco, Héctor, viéndolos, se arroja clamoroso,\\
con las recias falanges troyanas que comanda\\
Ares, al propio tiempo que Enío veneranda.\\
Va ésta con el Tumulto feroz y escandaloso,\\
mientras el dios, blandiendo su lanza prepotente,\\
ya marcha a espaldas de Héctor, ya anda con él al frente.\\
Tembló al verlo Diomedes el bravo, a la manera\\
del que yendo extraviado por una gran pradera,\\
dá con rápido río que echa al mar su corriente,\\
y al ver la hirviente espuma, salta retrocediendo,\\
así paró Tideides, a los suyos diciendo:\par
--Oh amigos ¿cómo habremos de asombrarnos que sea\\
Tan bravo el divino Héctor y audaz en la pelea?\\
Siempre con él va un numen que la muerte le evita,\\
Y ahora es Ares, que a un hombre mortal junto a él imita.
\pend
\separator
%
\pstart
\setline{720}
Aprontar los corceles de áurea brida, dispuso\\
la augusta diosa Hera, hija del gran Kronos, entonces.\\
A su vez la pronta Hebe fué y al coche le puso\\
en el eje de hierro combas ruedas de bronce,\\
de ocho rayos, que sobre pinas de oro inusable\\
llantas de bronce ostentan con ajuste admirable;\\
mientras que son de plata los cubos bien torneados.\\
Y de oro y plata tiene la caja sus sopandas,\\
y por la delantera le corren dos barandas.\\
El timón es de plata, y ella a su extremo ha atado\\
yugo y coyundas de oro de gran belleza. Y hera,\\
avida de combates y ejercicio agitado,\\
unce al yugo los potros firmes en la carrera.\par
Mientras tanto Atena, hija del dios portaéjida, echa\\
sobre el umbral paterno su velo bien bordado\\
que sus manos tejieron, y luego se ha ajustado\\
el arnés del nubígero dios: que así se pertrecha\\
para la guerra aciaga; y a la espalda ha cargado\\
la fiera y floqueada égida que el \emph{Espanto} corona.\\
En ella están la \emph{Fuerza}, la \emph{Discordia}, y la atroz\\
\emph{Persecución}, y la hórrida testa de la \emph{Gorgona},\\
prodigio de Zeus, monstruo tremebundo y feroz.\\
Armaste la cabeza con una áurea celada\\
de doble penacho y de cuádruple carrillera,\\
que a los infantes juntos de cien pueblos cubriera;\\
y montando el fulgente carro, ase la pesada\\
ingente y recia lanza con que a tanto valiente\\
rinde, si en contra suya se irrita prepotente.\par
\pend
\separator
\endnumbering